\chapter{Defining directional extensions} \label{ch:identifying}

This chapter defines and exemplifies the types of directional extensions found in Chadic languages. Section \ref{sec:path} further defines the morphosyntactic notion \textit{verbal extension} and the semantic concept \textit{direction}. The remaining sections each describe and exemplify a particular sub-type of directional extension found in Chadic languages: ventive (Section \ref{sec:vent1}), itive (Section \ref{sec:itive1}), inward and outward (Section \ref{sec:bounded1}), upward and downward (Section \ref{sec:vertical1}), `onto' and `under' (Section \ref{sec:onto1}). A few marginal types are presented in Section \ref{sec:otherDIR1}: homeward, bushward and `side'. The results of a quantitative study looking at how common each subtype of directional extensions is across the branches of Chadic languages is presented in Chapter \ref{ch:distribution}. Before that chapter, Chapter \ref{ch:identifying2} further explains how types of directional extensions were identified in published sources and organized into a database for quantitative analysis. 

\section{Directional extensions} \label{sec:path}

In the analysis of Chadic languages, verbal morphology is typically organized into three categories. First, there is normally verbal morphology expressing notions related to tense-aspect. Second, there is often a type of verbal morphology that has argument agreement or pronominal functions. Third, verbal morphology not falling into one of those two categories are frequently described as \textit{verbal extensions}.

Verbal extensions may be affixes or particles (morphologically detached from the verb) but all extensions are morphosyntactically dependent on the verb. Non-concatenative verbal morphology and irregular verbal conjugations are also considered types of verbal extensions. Unlike adverbs, a verbal extension cannot modify a nonverbal predicate or be used independently of a verb. Even if the verb may be elided (e.g., in response to a question) a verbal extension cannot be uttered on its own without a verb. Verbal extensions have fixed positions in the verbal phrase (typically following the verb), and do not exhibit variable syntactic positioning such as being able to appear at the beginning of a clause. In principle, the distinction between verbal extensions and other categories is not fuzzy, but it is not always clear how to diagnose a particular morpheme from the available documentation and description.

Practically speaking, available descriptions are generally taken at face value, and any directional morpheme described as a part of the verbal morphology is usually considered an extension. However, there are a few disputable descriptions. In \hyperref[tera]{Tera}, \citet[18]{newmanGrammarTeraTransformational1970} describes the form \textit{ɓara} as a ``particle'' but \citet[314--315]{schuhChadicCornucopia2017} points out that it has the syntactic distribution of an adverb. Some \hyperref[masa]{Masana} morphemes are described as verbal extensions, but also have uses in nonverbal clauses which (at least in this book) places them in the category of adverbs \citep[226]{melisDescriptionMasaTchad1999}. In a description of \hyperref[bura]{Bura}, \citet[284--289]{hoffmannUntersuchungenZurStruktur1955} describes several forms as phonetic contractions of a verb and an adposition or noun, rather than as productive verbal morphology. These few examples are not considered verbal extensions in this work.

One common function of verbal extensions in Chadic languages is to add directional information about a path of motion.\footnote{In addition to directional extensions, another common function of verbal extensions in Chadic languages is to modify the argument structure or valency of the verb.} Directional extensions are verbal extensions that can (as at least one of their uses) be added to a verb expressing a (translocative) motion event in order to specify the orientation of the path of motion. Directionals are frequently shown to combine with intransitive verbs of motion, for example, a manner-of-motion verb, as in (\ref{ex:muse1}), or a directional verb, as in (\ref{ex:bura6}). Note that in (\ref{ex:bura6}), information about the direction of the path of motion comes from both the ventive directional verb as well as the downward directional extension.  

\ea\label{ex:muse1}
\langinfo{\hyperref[muse]{Musey} verb \textit{līŋ} `run' with ventive marker}{}{\cite[76]{dassidiMorphophonologicalStructureMusej2015}} \\
\gll Sánàmā līŋ-áj tá ànú	\\
Sanama run-\gl{vent} body \gl{1}\gl{sg}.\gl{obj} \\
\glt `Sanama ran towards me.'

\ex\label{ex:bura6}
\langinfo{\hyperref[bura]{Bura} verb \textit{si} ‘come’ with the downward suffix}{}{\cite[281]{hoffmannUntersuchungenZurStruktur1955}} \\
\gll si-hi \\
come-\gl{down} \\
\glt `come down' [German: `\textit{herunterkommen}']
\z 

Semantically there are eight predominant types of directional meaning expressed in verbal extensions in Chadic languages. These can be seen as four pairs of semantic counterparts: ventive-itive, inward-outward, upward-downward and onto-under. Ventive and itive direction are defined in relation to a deictic center. Ventive is movement toward a deictic center (deictic center is the goal) (Section \ref{sec:vent1}), and itive is motion away from a deictic center (Section \ref{sec:itive1}).\footnote{Note that itive motion is only the antonym of ventive motion if the source of the itive path of motion is the deictic center. See Section \ref{sec:itive1} for discussion of the itive directional extension in \hyperref[buwa]{Buwal} which has a source of motion other than the location of the speaker.} The deictic center is a relative notion, like the adverbial `here'. An understanding of where the deictic center is understood to be can only be derived from context. 

Upward and downward motion are defined in relationship to vertical space, with the earth's surface (or the direction of gravity) as the point of reference (Section \ref{sec:vertical1}). Inward and outward direction are defined in relation to a bounded or enclosed space (as in `enter' and `exit') (Section \ref{sec:bounded1}). Directional extensions meaning `onto' describe a path of motion which moves toward the uppermost surface of some object (Section \ref{sec:onto1}). The counterpart is a more rare `under' directional extension that describes a path of motion ending at or passing through a point below some object. There are also a few even rarer types of directional extensions discussed in Section \ref{sec:otherDIR1}.

\section{Ventive directional extensions} \label{sec:vent1}

A ventive directional extension combined with a motion verb has the function of indicating that the path of motion is directed toward the deictic center. In many contexts, the deictic center will be understood to be the location of the speaker, but this will not always be the case. Precisely how the deictic center is pragmatically interpreted will not be explored here. It will simply be inferred from available translations whether a particular form expresses movement toward or away from a deictic center.\footnote{Assuming that translations consistently reflect the deictic orientation is not without complications. In the descriptions of four languages, there is at least one example of a ventive extension that is given an exceptional itive translation (\hyperref[mina]{Mina}, \hyperref[mbuk]{Mbuku}, \hyperref[mbud]{Mbudum}, \hyperref[mada]{Mada}), and in two examples an itive extension is given a ventive translation (\hyperref[zulg]{Gemzek} and \hyperref[gaan]{Ga'anda}). These exceptions are left aside without explanation.} 

Several different labels are used for ventive directional extensions in descriptions of Chadic languages, including \textit{centripetal} and, in French, \textit{rapprochement}. They are sometimes simply called \textit{directional}. A ventive extension may be found with a manner-of-motion verb to indicate in which direction the motion occurs, as in (\ref{ex:vame1}) and (\ref{ex:mofu1}).

\ea\label{ex:vame1}
\langinfo{\hyperref[vame]{Vamé} verb \textit{hav} `fly' with ventive suffix}{}{\cite[56]{kinnairdVameVerbalSystem2006}} \\
\gll  fan hav-ki-ke	\\
bird fly-\gl{vent}-\gl{prf}	\\
\glt `The bird is flying (towards us).'

\ex\label{ex:mofu1}
\langinfo{\hyperref[mofu]{Mofu} verb \textit{há} `run' with ventive suffix}{}{\cite[70]{barreteauDescriptionMofuGudurLangue1983}} \\
\gll há-wa \\
run-\gl{vent} \\
\glt `Come here quickly!' [French: `\textit{Viens-vite ici} !']
\z 

A ventive extension can also combine with a directional motion verb that indicates a different parameter of direction, such as the downward motion verb in (\ref{ex:tang1}) which combines with a ventive extension indicating that the motion is both downward and toward the deictic center (the inverse pattern of \ref{ex:bura6}). 

\ea\label{ex:tang1}
\langinfo{\hyperref[tang]{Tangale} verb \textit{yẹk} `descend' with ventive suffix}{}{\cite[47]{jungraithmayrDictionaryTangaleLanguage1991}} \\
\gll n yẹk-tụ́-ngọ	\\
\gl{1}\gl{sg} descend-\gl{vent}-\gl{1}\gl{sg}	\\
\glt `I came down.'	
\z

In several languages, there is a general motion verb which does not entail any specific direction of its path (Chapter \ref{ch:verbs}, Section \ref{sec:motionV}). In a neutral context, the implied path might be itive (away from the deictic center) or ventive. In the two instances of the \hyperref[buwa]{Buwal} verb \textit{nda} in (\ref{ex:buwa1}) each token has a different interpretation (in the original text, one is glossed `come' and one is glossed `go'). The same verb can combine with a ventive extension to explicitly express ventive motion, as in (\ref{ex:buwa100}). 

\ea\label{ex:buwa1}
\langinfo{\hyperref[buwa]{Buwal} verb \textit{nda} meaning `come' and `go' }{}{\cite[365]{viljoenGrammaticalDescriptionBuwal2013}} \\
\gll  màdā māwàl ká-\textbf{ndā} āzá nèné-ná-\textbf{ndā} á eɡljz \\
if husband \gl{pfv}-go/come \gl{compl} \gl{1}\gl{excl}.\gl{sbj}-\gl{fut}-go/come \gl{prep} church \\
\glt `If my husband has come, we will go to church.'

\ex\label{ex:buwa100}
\langinfo{\hyperref[buwa]{Buwal} verb \textit{nda} `go/come' with ventive extension}{}{\cite[373]{viljoenGrammaticalDescriptionBuwal2013}} \\
\gll ná-\textbf{ndā}-xā ná-ndzā á bwāl \\
\gl{1}\gl{excl}.\gl{sbj}-go/come-\gl{vent} \gl{sbj}.\gl{1}\gl{pl}.\gl{excl}-stay \gl{prep} Buwal	\\
\glt `...we came, we stayed in Buwal village. (Speaker is located in Buwal village.)'
\z

Ventive extensions can also be used with transitive (caused) motion verbs where the object is a figure on a path of motion. In (\ref{ex:shaa2}), the translation describes an event in which both the subject and object are moving along a path of motion toward the deictic center.

\ea\label{ex:shaa2}
\langinfo{\hyperref[shaa]{Ron-Scha} verb \textit{də̀ŋ} `pull' with ventive suffix}{}{\cite[268]{jungraithmayrRonSprachenTaschadohamitischeStudien1970}} \\
\gll də̀ŋ-ó \\
pull-\gl{vent} \\
\glt `drag here' [German: `\textit{her-ziehen}']
\z 

\section{Itive directional extensions} \label{sec:itive1}

Itive direction is the counterpart of ventive, describing a path of motion that leads to a location farther from the deictic center. For example, in (\ref{ex:buwa10}), the \hyperref[buwa]{Buwal} itive extension \textit{āzà} combines with a manner-of-motion verb to express that the figure on the path of motion moves away. 

\ea\label{ex:buwa10}
\langinfo{\hyperref[buwa]{Buwal} verb \textit{xēj} `run' with itive extension}{}{\cite[280]{viljoenGrammaticalDescriptionBuwal2013}} \\
\gll ā-xēj āntā āzà skʷá	\\
\gl{3}\gl{sg}.\gl{sbj}-ran \gl{3}\gl{sg}.\gl{poss} \gl{itive} \gl{q}	\\
\glt `…did he run away?'	
\z

Itive directionals are generally presented as antonyms of ventive directionals, but this is not always the case. In \hyperref[buwa]{Buwal}, the ventive and itive directional extensions can co-occur, as in (\ref{ex:buwa01}). This is possible because they do not have the same deictic reference point.

\ea\label{ex:buwa01}
\langinfo{\hyperref[buwa]{Buwal} verb \textit{tàdàkʷ} `descend' with ventive and itive directionals}{}{\cite[378]{viljoenGrammaticalDescriptionBuwal2013}} \\
\gll xʷné-tàdàkʷ-xā āzà á tā xājāk \\
2\gl{pl}.\gl{sbj}-descend-\gl{vent} \gl{itive} \gl{prep} on ground\\ 
\glt `You come down from there onto the ground! (Speaker on the ground.)'
\z

According to \citet[378]{viljoenGrammaticalDescriptionBuwal2013}, the source of the `away' itive motion cannot be the speaker. In (\ref{ex:buwa01}), the source location is up in a tree while the speaker is on the ground. The goal of the ventive motion is the speaker, who is on the ground. The combination of ventive and itive can be analyzed as `move towards me (on the ground) [ventive] and away from that tree [itive]'. Other descriptions of Chadic languages do not necessarily contain this level of detail about the source of an itive path of motion. For this study, all directional extensions described as expressing that a path of motion is `away' are grouped together as itive irrespective of how the source of that path is defined.

As mentioned, some Chadic languages have a general motion verb without any inherent directional meaning, such as \textit{d} in \hyperref[cuvo]{Cuvok} which can be combined either with a ventive extension, as in (\ref{ex:cuvo1}), or an itive extension, as in (\ref{ex:cuvo2}).

\ea\label{ex:cuvo1}
\langinfo{\hyperref[cuvo]{Cuvok} verb \textit{d} `go/come' with ventive suffix}{}{\cite[303]{dadakGrammaireCuvokLangue2021}} \\
\gll d-ɛ́k	\\
go/come-\gl{vent} \\
\glt `Come here.'

\ex\label{ex:cuvo2}
\langinfo{\hyperref[cuvo]{Cuvok} verb \textit{d} `go/come' with itive suffix}{}{\cite[308]{dadakGrammaireCuvokLangue2021}} \\
\gll d-áɗ \\
go/come-\gl{itive} \\
\glt `Go over there.'
\z 

Itive extensions are also found with transitive motion verbs in which the patient is the figure on a path of motion away from the deictic center, as in (\ref{ex:marg3}). 

\ea\label{ex:marg3}
\langinfo{\hyperref[marg]{Margi} verb \textit{ndàl} `throw' with itive suffix}{}{\cite[131]{hoffmannGrammarMargiLanguage1963}} \\
\gll ndàl-nà	\\
throw-\gl{itive}		\\
\glt `to throw away'
\z 

\section{Boundary-crossing directional extensions} \label{sec:bounded1}

Boundary-crossing directional extensions express one of a pair of directional meanings: inward or outward. Inward and outward directional extensions can be combined with directional motion verbs, such as the ventive and itive verbs in (\ref{ex:psik12}) and (\ref{ex:saku5}). They can also be combined with manner-of-motion verbs, as in (\ref{ex:marg6b}).

\ea\label{ex:psik12}
\langinfo{\hyperref[psik]{Psikye} verb \textit{dze} `go' with inward suffix}{}{\cite[118]{smithKapsikiLanguage1969}} \\
\gll pa dze-mbé cɛ	\\
then go-\gl{into} house	\\
\glt `Then (he) went into the house...'

\ex\label{ex:saku5}
\langinfo{\hyperref[saku]{Sakun} verb \textit{ja} `come' with outward suffix}{}{\cite[126]{thomasGrammarSakunSukur2014}} \\
\gll má ja-vá-j ja-vá kwá ná … \\
when come-\gl{out}-\gl{prog} come-\gl{out} \gl{2}\gl{sg} \gl{top} \\
\glt `Immediately when you come out...'

\ex\label{ex:marg6b}
\langinfo{\hyperref[marg]{Margi} verb \textit{mə̀} `run' with outward suffix}{}{\cite[122]{hoffmannGrammarMargiLanguage1963}} \\
\gll mə̀-bá	\\
run-\gl{out}		\\
\glt `to start up and run out'
\z

Inward and outward extensions are also found with transitive motion events, as in (\ref{ex:marg12}) and (\ref{ex:bura16}).

\ea\label{ex:marg12}
\langinfo{\hyperref[marg]{Margi} verb \textit{f} `put' with inward suffix}{}{\cite[148]{hoffmannGrammarMargiLanguage1963}} \\
\gll  f-wá	\\
put-\gl{into} \\
\glt `to put (one) into'

\ex\label{ex:bura16}
\langinfo{\hyperref[bura]{Bura} verb \textit{buka} ‘push’ with outward suffix}{}{\cite[4]{blenchBuraVerbalExtensions2010}} \\
\gll  buka-ɓəla \\
push-\gl{out} \\
\glt `to push outside'
\z  

Inward and outward directional meanings are defined in relation to any space perceived as having an enclosed boundary. This can be as concretely defined as a jar or a house, as in (\ref{ex:psik12}), or it can be more abstractly defined, for example, as a fire which has no fixed borders, as in (\ref{ex:dghw800}).

\ea\label{ex:dghw800}
\langinfo{\hyperref[dghw]{Dghwede} verb \textit{ə̀yə̂} `fall' with inward suffix}{}{\cite[25]{frickVerbalSystemDghwede1978}} \\
\gll ə̀yə̂-me cé ne dá me kárà... \\
fall-\gl{into} \gl{past} ?? to ?? fire \\
\glt `when he had fallen into the fire'
\z 

In \hyperref[mata]{Matal}, there are two inward directional extensions which are semantically distinguished by whether the bounded area is enclosed (e.g., with a cover) or not. The form \textit{ába} is used for uncovered containers, such as the pipe in (\ref{ex:mata10}), and \textit{águ} is used for covered containers, as in (\ref{ex:mata11}) \citep[76]{dougopheGrammarMatal2021}.\footnote{Note that further data is needed to more definitively determine the morphosyntactic status of these morphemes---whether they are part of the verbal system or adverbials.}

\ea\label{ex:mata10}
\langinfo{\hyperref[mata]{Matal} verb \textit{d} `go' with inward extension}{}{\cite[180]{dougopheGrammarMatal2021}} \\
\gll áŋá ákàl à-dá-d ábà jkə̀nj á-s-àl=àw	\\
if fire \gl{3}\gl{sg}-\gl{pfv}-go inside too \gl{3}\gl{sg}-please-\gl{3}\gl{sg}.\gl{dat}=\gl{neg}	\\
\glt `If heat enters the pipe, it does not like/it is not good.'

\ex\label{ex:mata11}
\langinfo{\hyperref[mata]{Matal} verb \textit{sàkw} `put' with inward extension}{}{\cite[202]{dougopheGrammarMatal2021}} \\
\gll  wànáj xə̀ŋ ka mə̀-sàkw-àx-àj ka kwàndúrkú águ	\\
\gl{dem} (\gl{prox}) \gl{top} \gl{1}\gl{sg}-put-\gl{pl}-?? \gl{top} dolo inside		\\
\glt `As for this one, we put dolo (an unfermented sorghum beverage) inside.'
\z 

\section{Vertical directional extensions} \label{sec:vertical1}

Vertical directional extensions express upward or downward direction. Upward direction can be thought of as a path of motion that moves counter to the direction of gravity, and downward direction is a path of motion in which the figure moves in the same direction as gravity. Upward and downward extensions can combine with directional motion verbs, as in (\ref{ex:malg1}) and (\ref{ex:lama8}), and also with manner-of-motion verbs, as in (\ref{ex:lama15}). 

\ea\label{ex:malg1}
\langinfo{\hyperref[malg]{Malgwa} verb \textit{d} `go' with upward suffix }{}{\cite[153]{lohrSpracheMalgwaNara2002}} \\
\gll də́-tə́ \\
go-\gl{up} \\
\glt `to climb (e.g., up a tree)'

\ex\label{ex:lama8}
\langinfo{\hyperref[lama]{Lamang} verb \textit{s} ‘come’ with downward suffix}{}{\cite[144]{wolffLamangLanguageDictionary2015}} \\
\gll sa-gáa-tá	\\
come-\gl{down}-\gl{compl}	 \\
\glt `coming down'

\ex\label{ex:lama15}
\langinfo{\hyperref[lama]{Lamang} verb \textit{ndr} ‘fly’ with downward suffix}{}{\cite[283]{wolffLamangLanguageDictionary2015}} \\
\gll ɗíyák ndə̀rà-gáa-tá skwè-ɗè	\\
bird fly-\gl{down}-\gl{compl} come-\gl{down}		\\
\glt `The bird flew down.'
\z 

Like other directional extensions, vertical directional extensions are also used with transitive motion verbs, as in (\ref{ex:bura17}). 

\ea\label{ex:bura17}
\langinfo{\hyperref[bura]{Bura} verb \textit{ndzi} ‘throw’ with downward suffix}{}{\cite[281]{hoffmannUntersuchungenZurStruktur1955}} \\
\gll ndzi-hi \\
throw-\gl{down} \\
\glt `throw down' [German: `\textit{niederwerfen}'] 
\z 

\citet[141]{wolffLamangLanguageDictionary2015} describes the \hyperref[lama]{Lamang} upward extension \textit{-fi} as also expressing ``direction east'' and the downward extension \textit{-ɗi} as also expressing ``direction west''. This corresponds to Lamang speakers residing in plains to the west of the northernmost tip of the Mandara mountain range in eastern Nigeria. From this vantage point, moving upwards into the mountains is also moving eastward. Conversely, in \hyperref[glav]{Glavda}, \citet[662]{bubaGlavdaMorphology2007} describe the upward extension \textit{-dit} as having an additional meaning of ``westward'', as in (\ref{ex:glav7}). This corresponds to Glavda speakers residing in an area to the east of the same mountains, and so from their perspective movement into the mountains is movement westward.\footnote{However, \citet{nghagyiyaSketchGlavdaGrammar2011} contradicts \citet[662]{bubaGlavdaMorphology2007} and instead associates the upward extension with `eastward'.} 

\ea\label{ex:glav7}
\langinfo{\hyperref[glav]{Glavda} verb \textit{ɬəg} `push' with \textit{-dít}  suffix }{}{\cite[663]{bubaGlavdaMorphology2007}} \\
\gll ɬəg-a-d-ít-ɬəga \\
push-\gl{3}-\gl{tr}-\gl{up}-push \\
\glt `He pushed something up/from east to west.'
\z

Similar patterns of using geographic landmarks to associate cardinal direction with a relative deictic direction have been observed in several languages around the world. For example, in varieties of the Maasai language (Nilotic) words expressing `top' and `bottom' can co-express the cardinal notions `north' and `south' depending on the geographic position of the homeland of the speakers in relation to Mount Kilimanjaro \citep{mietznerSpatialOrientationNilotic2012}. \citet{barthDirectionalConstructionsMatukar2015} describe directional morphemes in Matukar Panau (Oceanic) in which `upward' meaning is co-expressed with `inland' while `downward' meaning is co-expressed with `seaward'. 

In Chadic languages, it is not clear how the cardinal directional meaning functions for speakers who travel away from these particular areas. Would the cardinal meaning be retained for speakers who traverse from one side of the mountain range to the other? Or is it only in a particular locality that upward movement and downward movement are extended to mean eastward and westward? 

\section{Directional extensions `onto' and `under'}  \label{sec:onto1}

The direction `onto' refers to a path of motion which ends at the uppermost surface of a particular object.\footnote{There is one example in the \hyperref[malg]{Malgwa} data where the `onto' extension \textit{ar} is given an `over' translation with the verb \textit{bə́z} `jump' \citep[159]{lohrSpracheMalgwaNara2002}. The same translation is also found with the upward extension \textit{tə́} and the same verb \citep[162]{lohrSpracheMalgwaNara2002}. In the database, these few instances of an `onto' extension with an `over' translation are included as examples of an `onto' extension.}  These most frequently appear in examples with transitive motion verbs, as in (\ref{ex:dghw13}), but also appear with intransitive motion verbs, as in (\ref{ex:bura8}).  

\ea\label{ex:dghw13}
\langinfo{\hyperref[dghw]{Dghwede} verb \textit{pàɗ} `pour' with `onto' suffix}{}{\cite[26]{frickVerbalSystemDghwede1978}} \\
\gll ká' tè pàɗ-àrə̀-gà páɗá yíwè nè kwíré yíwè \\
\gl{narr} ?? pour-\gl{dat}-\gl{onto} pour water ?? stone water	\\
\glt `They were repeatedly pouring water on the water stone.'

\ex\label{ex:bura8}
\langinfo{\hyperref[bura]{Bura} verb \textit{ndla} ‘fall’ with `onto' suffix}{}{\cite[7]{blenchBuraVerbalExtensions2010}} \\
\gll ndla-nkər \\
fall-\gl{onto} \\
\glt `to fall on top of something'
\z 

Three languages have an `under' directional extension which describes a path of motion toward the bottom side of object. There are only seven examples of these directional extensions in the data, one of which is with an intransitive motion verb, shown in (\ref{ex:park8}). The remaining examples are found with transitive motion verbs, as in (\ref{ex:glav8}).

\ea\label{ex:park8}
\langinfo{\hyperref[park]{Parkwa} verb \textit{d} `go' with `under' suffix}{}{\cite[315]{jarvisPodokoVerbalDirectionals1983}} \\
\gll D-asə mayə́ akə́ sla də́ zlə́ma \\
go-\gl{under} \gl{1}\gl{sg} \gl{prep} cow \gl{prep} stable	\\
\glt `I went into the cowshed. (Literally, I went under to the cow in the stable. The stable is lower than the rest of the house.)'

\ex\label{ex:glav8}
\langinfo{\hyperref[glav]{Glavda} verb \textit{f} `put' with `under' suffix }{}{\cite[24]{nghagyiyaSketchGlavdaGrammar2011}} \\
\gll f-àn-ár-s	fᵊ-g	k-ká:rá \\
put-\gl{1}\gl{sg}-??-\gl{under}	put-??	\gl{obj}-fire \\
\glt `I put fire under it.'
\z

\section{Homeward and bushward directional extensions} \label{sec:otherDIR1} \label{sec:home} \label{sec:bush}

The most commonly attested directional extensions are discussed in the preceding sections of this chapter. This section briefly describes a number of marginal examples of directional extensions that are only attested in a few examples in a few languages. A homeward directional extension is described in the grammar of two languages, and one language has a bushward directional extension, although it is only attested in this function with the itive motion verb `go'. 

The Biu-Mandara languages \hyperref[lama]{Lamang} and \hyperref[bura]{Bura} have a directional extension which describes a path of motion in a homeward direction, and two other languages (\hyperref[psik]{Psikye} and \hyperref[huba]{Huba}) have a nascent homeward morpheme which only occurs with the basic motion verbs (Chapter \ref{ch:diachrony}, Section \ref{sec:diachronyother}).\footnote{\citet{smithMuyangVerbPhrase2002} describes a verbal suffix \textit{-biyu} in \hyperref[muya]{Muyang} as `homeward' but its meaning is ventive, and in a subsequent publication \citet{smithDefinitenessTopicalisationTheme2003} labels the same suffix `hither'.} In his description of Lamang, Wolff uses the label ``ventive'' for the homeward extension, as in (\ref{ex:lama3}), but clarifies that the  meaning is ``movement into and/or towards an inhabited or sacred place'' \citep[140]{wolffLamangLanguageDictionary2015} and elsewhere states that the form is neither ventive nor itive \citep[234]{wolffEncodingTopographyDirection2006}. There are two homeward directional forms in Lamang. The suffix \textit{-gha} is the form of the homeward extension with the verb \textit{dza-} `go' and \textit{skwa-} `come'.\footnote{\citet[140]{wolffLamangLanguageDictionary2015} identifies this form as derived from the noun \textit{tə́ghà} `compound, residence, home'.} With other verbs, the homeward form is marked by vowel lengthening and a high-low tone pattern \citep[151]{wolffLamangLanguageDictionary2015}. Note that in (\ref{ex:lama3}) the homeward extension is used in an itive motion context with the verb \textit{dzá} `go' which shows that it is not expressing the contradictory ventive meaning.

\ea\label{ex:lama3}
\langinfo{\hyperref[lama]{Lamang} verb \textit{dzá} ‘go’ with homeward suffix}{}{\cite[211]{wolffLamangLanguageDictionary2015}} \\
\gll gú plís dzá-ghà dáa Wándàlà \\
\gl{narr} horse go-\gl{home} \textsc{prep} Wandala \\
\glt `And then the horse(man) went home to Wandala.'
\z 

One example of a transitive motion verb with a homeward suffix is found in \hyperref[bura]{Bura}, as shown in (\ref{ex:bura10}).\footnote{Blench's list of verbal extensions also includes the suffix \textit{-ha} which is translated as `to an inhabited place, home' but is elsewhere described as meaning `together, collectively'.}  \citet{hoffmannUntersuchungenZurStruktur1955} takes examples from a Bura translation of the New Testament. The context of (\ref{ex:bura10}) is a story in which the friends of a sick man carry his stretcher onto the roof of a crowded house and remove part of the roof in order to lower the friend into the house on the stretcher, thereby gaining access to Jesus who miraculously heals the man.

\ea\label{ex:bura10}
\langinfo{\hyperref[bura]{Bura} verb \textit{ha} ‘put’ with the homeward suffix}{}{\cite[301]{hoffmannUntersuchungenZurStruktur1955}} \\
\gll ha-vi \\ 
put-\gl{home} \\
\glt `let down (into the house)' [German: `\textit{hinunterlassen} (\textit{scil. ins Haus})'] 
\z

As a semantic parallel to homeward direction in verbal morphology, it is worth mentioning that the \hyperref[marg]{Margi} lexicon includes two motion verbs which can describe inward or homeward motion: \textit{lì}  `go in, go (home)' and \textit{sì}  `come in, come (home)'.

In African contexts, \textit{the bush} refers to uninhabited areas, in contrast with villages, towns and cities. In \hyperref[psik]{Psikye}, \citet[118]{smithKapsikiLanguage1969} describes a suffix \textit{-mté} which can combine with the motion verb \textit{dze} `go' to mean `go to the bush', as in (\ref{ex:psik13}). 

\ea\label{ex:psik13}
\langinfo{\hyperref[psik]{Psikye} verb \textit{dze} `go' with bushward suffix}{}{\cite[118]{smithKapsikiLanguage1969}} \\
\gll ntíŋú'yá ka-dze-mté \\
suddenly \gl{1}\gl{sg}-\gl{ipfv}-go-\gl{bush}	\\
\glt `I abruptly went to the bush.'
\z 

However, the bushward directional meaning of this verbal suffix is only attested with the verb \textit{dze} `go', and even with this verb, some translations give a more general itive interpretation, as in (\ref{ex:psik100}).

\ea\label{ex:psik100}
\langinfo{\hyperref[psik]{Psikye} verb \textit{dze} `go' with buswhward suffix}{}{\cite[118]{smithKapsikiLanguage1969}} \\
\gll dze-mté jaxaŋgalé zaaŋkɛ́	\\
go-\gl{bush} trip her.man	\\
\glt `Her husband has left on a trip.'
\z 

This chapter has presented the semantic types of directional extensions found in Chadic languages. Chapter \ref{ch:identifying2} provides additional details of how a database suitable for quantitative analysis was constructed on the basis of published sources, and then Chapter \ref{ch:distribution} presents a quantitative analysis of how many Chadic languages in each branch have evidence of various types of directional extensions. Note that since homeward and bushward directional extensions are so rare, they will not be included in the quantitative analysis. 
